% Bonus: exercicios classicos de programacao inteira por nivel (GCC118, UFLA)
% Compilar a partir desta pasta:  tectonic desafios_bonus_classicos.tex   (ou pdflatex/xelatex, duas vezes)
\documentclass[11pt,a4paper]{article}

\usepackage{iftex}
\ifPDFTeX
  \usepackage[utf8]{inputenc}
  \usepackage[T1]{fontenc}
  \usepackage{lmodern}
\else
  \usepackage{fontspec}
\fi
\usepackage[brazil]{babel}
\usepackage[margin=2.5cm]{geometry}
\usepackage{amsmath,amssymb}
\usepackage{icomma}
\usepackage{booktabs,array}
\usepackage{enumitem}
\usepackage[hidelinks]{hyperref}

\setlist{leftmargin=*,itemsep=1pt,topsep=3pt}
\setlength{\parindent}{0pt}
\setlength{\parskip}{5pt}
\setlength{\emergencystretch}{3em}
\renewcommand{\arraystretch}{1.15}
\allowdisplaybreaks

% Bloco numerado: "Bloco 1 -- Roteamento"; subsecoes ficam 1.1, 1.2, ...
\newcommand{\bloco}[1]{%
  \clearpage\refstepcounter{section}%
  \section*{Bloco \thesection\ --- #1}%
  \addcontentsline{toc}{section}{Bloco \thesection\ --- #1}}
\newcommand{\enunciado}{\par\textbf{Enunciado.}\ }
\newcommand{\tarefa}{\par\textbf{Tarefa.}\ }
\newcommand{\passo}[1]{\par\medskip\textbf{#1}\par}
\newcommand{\porque}{\textit{Por quê?}\ }
\newenvironment{nota}[1]{\par\smallskip\begin{list}{}{\leftmargin=1.2em\rightmargin=0pt}\item[]\textbf{#1.}\ }{\end{list}}

\title{\textbf{Bônus: exercícios clássicos de Programação Inteira}\\[4pt]\large Por nível de dificuldade, com resolução comentada}
\author{GCC118 --- Programação Matemática (UFLA)\\[2pt]\small Complemento ao documento \emph{Desafios de Programação Inteira --- resolução comentada}}
\date{}

\begin{document}
\maketitle
\thispagestyle{empty}

\begin{abstract}
\noindent Este documento bônus complementa a resolução comentada dos desafios. Ele começa explicando com calma o que é o
big-M, a técnica que aparece em quase todos os modelos, e depois traz onze exercícios de problemas clássicos de
programação inteira, nos moldes dos apresentados em Arenales et al.\ (2007), organizados em três níveis de dificuldade.
Todos os valores ótimos citados foram conferidos com PuLP/CBC.
\end{abstract}

\tableofcontents

\clearpage
\section*{O que é o big-M (restrição que liga e desliga)}
\addcontentsline{toc}{section}{O que é o big-M}

\textbf{O problema que ele resolve.} Num modelo de programação linear ou inteira, toda restrição é do tipo
``combinação linear das variáveis $\le$ (ou $\ge$, $=$) número'', e vale \emph{sempre}. Não existe a palavra ``se''. Mas os
enunciados estão cheios de condições: \emph{se} houver pernoite, pode rodar mais de 300\,km; \emph{se} a máquina não for
alugada, não se produz nada; o pedido termina antes da manutenção \emph{ou} começa depois dela. O big-M é o truque que
simula esse ``se'' usando uma variável binária e uma constante grande, chamada $M$.

\textbf{A ideia.} Some ao lado direito da restrição um termo que vale 0 quando a restrição deve valer e vale $M$, um
número grande, quando ela deve ser ignorada. Exemplo, tirado do exercício 1.1 dos Desafios: um motorista que roda mais de
300\,km no dia não volta à base e a empresa paga um pernoite. Com $x$ = km rodados e $y=1$ se há pernoite, a regra
``se $y=0$ (não há pernoite), então $x\le300$'' fica:
\[
  x \le 300 + M\,y .
\]
\begin{center}
\begin{tabular}{ccl}
\toprule
$y$ & a restrição vira & efeito\\
\midrule
0 & $x \le 300$ & \textbf{ativa}: a regra vale\\
1 & $x \le 300 + M$ & \textbf{desligada}: com $M$ grande, nunca atrapalha\\
\bottomrule
\end{tabular}
\end{center}
Uma única restrição linear passa a se comportar de dois jeitos, conforme o valor da binária. Essa é a ideia inteira do big-M.

\textbf{As formas gerais.} Para uma restrição $g(x)\le b$ (lado esquerdo $g(x)$, lado direito $b$):
\begin{center}
\begin{tabular}{ll}
\toprule
Quero que a restrição valha quando\ldots & Escrevo\\
\midrule
$y=1$ & $g(x) \le b + M\,(1-y)$\\
$y=0$ & $g(x) \le b + M\,y$\\
$y=1$, restrição do tipo $\ge$ & $g(x) \ge b - M\,(1-y)$\\
\bottomrule
\end{tabular}
\end{center}
Dica: o termo com $M$ deve ser \emph{zero} exatamente quando a restrição deve valer. Se ficar em dúvida, teste
$y=0$ e $y=1$ e veja o que sobra.

\textbf{Os três usos mais comuns.}
\begin{enumerate}
  \item \textbf{Ligar uma quantidade a uma decisão sim/não:} $x\le M\,y$, ou seja, ``se $y=0$, então $x=0$''.
        Aparece em custo fixo, setup e abertura de depósito (exercício I1 deste bônus; 5.2 e 5.3 dos Desafios). Aqui o
        $M$ é o maior valor que $x$ pode assumir.
  \item \textbf{Ativar ou desativar uma restrição:} como no exemplo do pernoite e no encadeamento de tempos com janelas
        (1.3 dos Desafios).
  \item \textbf{``Ou-ou'' (disjunção):} duas restrições, uma com $M\,y$ e outra com $M\,(1-y)$. Com $y=0$ só a primeira
        vale; com $y=1$, só a segunda (I2 e A2 deste bônus; 3.1 dos Desafios).
\end{enumerate}

\textbf{Como calcular o $M$ a partir dos dados.} Siga três passos:
\begin{enumerate}
  \item escreva como a restrição fica quando está \emph{desligada}: $g(x)\le b+M$;
  \item descubra o \textbf{maior valor} que $g(x)$ pode assumir em qualquer solução que deva continuar permitida, usando os
        limites do problema (capacidades, horizonte de tempo, demandas);
  \item tome $M = (\text{maior valor de } g(x)) - b$.
\end{enumerate}
No pernoite, o maior total de km que cabe na jornada de 8\,h é 509,375 (dado do exercício), então $M=509{,}375-300\approx210$. Na ligação $x\le My$, o maior valor
de $x$ é o próprio $M$: a capacidade, ou algo menor que os dados garantam.

\textbf{Por que não usar simplesmente $M=10^9$?} Porque, embora o modelo continue correto no papel, ele fica
\emph{pior} na prática, por dois motivos:
\begin{itemize}
  \item \textbf{Relaxação fraca.} O branch-and-bound resolve versões em que $y$ pode ser fracionário. Com $x\le My$,
        basta $y=x/M$. No pernoite, rodando 497,5\,km, a relaxação precisa de $y\ge197{,}5/M$:
        \begin{center}
        \begin{tabular}{ccc}
        \toprule
        $M$ & menor $y$ na relaxação & pernoite ``pago'' ($250\,y$)\\
        \midrule
        210 & 0,94 & R\$\,235\\
        420 & 0,47 & R\$\,118\\
        $10^6$ & 0,0002 & R\$\,0,05\\
        \bottomrule
        \end{tabular}
        \end{center}
        Com $M$ enorme, o custo fixo praticamente some da relaxação, e o limitante fica longe do ótimo. Como o
        branch-and-bound usa esse limitante para descartar ramos, ele poda menos e demora mais.
  \item \textbf{Erros numéricos.} Solvers aceitam uma binária ``quase zero'', como $y=10^{-7}$, por tolerância. Com
        $M=10^9$, esse $y$ ``zero'' libera $x\le100$, e a regra é violada sem aviso.
\end{itemize}
\textbf{Resumo:} o $M$ deve ser \emph{grande o bastante} para não cortar nenhuma solução válida e \emph{o menor possível}
dentro disso. Sempre que usar um big-M, indique de onde ele veio.

\clearpage
\section*{Como usar este bônus}
\addcontentsline{toc}{section}{Como usar este bônus}

Os exercícios a seguir cobrem os problemas clássicos de otimização discreta apresentados em Arenales et al.\ (2007):
mochila, designação, cobertura, escala de pessoal, custo fixo, restrições ``ou-ou'' e ``se-então'', corte, funções
lineares por partes, particionamento, job shop e roteamento. Os enunciados e os dados foram elaborados para este
material, no estilo dos exemplos do livro, e todos os ótimos foram conferidos com PuLP/CBC.

\begin{itemize}
  \item \textbf{Nível 1 --- Básico:} uma ideia de modelagem por exercício, só com binárias ou inteiras simples.
  \item \textbf{Nível 2 --- Intermediário:} exige uma técnica de ligação (big-M, ou-ou, se-então, padrões, segmentos).
  \item \textbf{Nível 3 --- Avançado:} modelos com várias famílias de restrições interagindo, próximos de problemas reais.
\end{itemize}
Sugestão de estudo: tente modelar sozinho a partir do enunciado e só depois leia o passo a passo.

\newcommand{\exercicio}[1]{\subsection*{#1}\addcontentsline{toc}{subsection}{#1}}

% =====================================================================
\clearpage
\section*{Nível 1 --- Básico}
\addcontentsline{toc}{section}{Nível 1 --- Básico}

\exercicio{B1. Seleção de projetos (mochila com regras lógicas)}
\enunciado Uma prefeitura tem R\$\,120 mil para investir e seis projetos candidatos, cada um com custo e retorno social
estimado (em R\$ mil). Cada projeto é feito inteiro ou não é feito. Regras: (a) os projetos 2 e 3 disputam o mesmo terreno,
então no máximo um deles pode ser feito; (b) o projeto 6 depende da estrutura do projeto 4; (c) os projetos 1 e 5 fazem
parte do mesmo convênio e devem ser feitos juntos ou nenhum deles. Maximize o retorno total.
\begin{center}
\begin{tabular}{lcccccc}
\toprule
Projeto & 1 & 2 & 3 & 4 & 5 & 6\\
\midrule
Custo $c_i$ & 40 & 30 & 25 & 50 & 20 & 35\\
Retorno $r_i$ & 60 & 45 & 35 & 80 & 25 & 50\\
\bottomrule
\end{tabular}
\end{center}

\passo{Passo 1 --- Variáveis}
$x_i\in\{0,1\}$: $x_i=1$ se o projeto $i$ é executado. \porque ``Feito inteiro ou não feito'' é uma escolha sim/não.

\passo{Passo 2 --- Modelo}
\begin{align*}
  \max\ & 60x_1 + 45x_2 + 35x_3 + 80x_4 + 25x_5 + 50x_6\\
  \text{s.a.}\ & 40x_1 + 30x_2 + 25x_3 + 50x_4 + 20x_5 + 35x_6 \le 120 && \text{(orçamento: a ``mochila'')}\\
  & x_2 + x_3 \le 1 && \text{(a) no máximo um}\\
  & x_6 \le x_4 && \text{(b) se 6, então 4}\\
  & x_1 = x_5 && \text{(c) os dois ou nenhum}
\end{align*}
\porque Cada regra lógica vira uma relação simples entre binárias. ``No máximo um'' é soma $\le1$. ``Se 6 então 4'' só
proíbe $(x_6,x_4)=(1,0)$, por isso usa $\le$. ``Juntos ou nenhum'' é a equivalência, que se escreve com igualdade.

\passo{Passo 3 --- Solução}
Ótimo: \textbf{175}, com os projetos 2, 4 e 6 (custo $30+50+35=115$). Sem as regras o ótimo seria 185 (projetos 1, 2 e 4),
mas esse plano faz o projeto 1 sem o 5 e viola a regra (c).

A relaxação linear vale $182{,}08$: ela usa $x_1=x_5\approx0{,}083$ para gastar os R\$\,5 mil que sobram do orçamento.
É o comportamento típico da mochila relaxada, que ``preenche'' a capacidade com um pedaço de item.

\begin{nota}{Erro típico}
Escrever ``se 6 então 4'' como $x_6=x_4$. Isso obrigaria a fazer o 6 sempre que o 4 fosse feito, o que a regra não pede.
\end{nota}

% ---------------------------------------------------------------------
\exercicio{B2. Designação de tarefas}
\enunciado Quatro funcionários (A, B, C, D) devem receber quatro tarefas, uma para cada. O tempo (em horas) que cada
funcionário leva em cada tarefa está na tabela. Minimize o tempo total.
\begin{center}
\begin{tabular}{c|cccc}
\toprule
 & T1 & T2 & T3 & T4\\
\midrule
A & 9 & 2 & 7 & 8\\
B & 6 & 4 & 3 & 7\\
C & 5 & 8 & 1 & 8\\
D & 7 & 6 & 9 & 4\\
\bottomrule
\end{tabular}
\end{center}

\passo{Passo 1 --- Variáveis e modelo}
$x_{ij}\in\{0,1\}$: o funcionário $i$ faz a tarefa $j$.
\begin{align*}
  \min\ & \sum_{i,j} t_{ij}\,x_{ij}\\
  \text{s.a.}\ & \sum_j x_{ij} = 1 \quad\forall i && \text{(cada funcionário, uma tarefa)}\\
  & \sum_i x_{ij} = 1 \quad\forall j && \text{(cada tarefa, um funcionário)}
\end{align*}

\passo{Passo 2 --- Solução}
Ótimo: \textbf{13}, com A$\to$T2 (2), B$\to$T1 (6), C$\to$T3 (1) e D$\to$T4 (4).

\passo{Passo 3 --- A lição: a relaxação já é inteira}
Resolvendo o modelo com $0\le x_{ij}\le1$ \emph{contínuas}, o solver devolve a mesma solução inteira, com valor 13.
\porque A matriz de restrições da designação é \emph{totalmente unimodular}: todo vértice do poliedro da relaxação tem
coordenadas inteiras. Na prática, a designação (assim como transporte e fluxo em redes com dados inteiros) pode ser
resolvida como programação linear, sem branch-and-bound. Isso não vale para os demais exercícios deste material.

% ---------------------------------------------------------------------
\exercicio{B3. Cobertura de conjuntos: bases de ambulância}
\enunciado Uma cidade tem 8 bairros e 6 locais candidatos para bases de ambulância. Um bairro está coberto se alguma
base instalada chega até ele em no máximo 10 minutos. A tabela mostra quais bairros cada local cobre. Todos os bairros
precisam ser cobertos. (a) Qual é o menor número de bases? (b) E se cada local tiver um custo de instalação
($50, 40, 45, 30, 55, 35$ para L1 a L6)?
\begin{center}
\begin{tabular}{lcccccc}
\toprule
Local & L1 & L2 & L3 & L4 & L5 & L6\\
\midrule
Bairros cobertos & 1, 2, 3 & 2, 4, 5 & 3, 5, 6 & 4, 7 & 6, 7, 8 & 1, 8\\
\bottomrule
\end{tabular}
\end{center}

\passo{Passo 1 --- Variáveis e modelo}
$y_\ell\in\{0,1\}$: instalar base no local $\ell$. Para cada bairro, escreva quais locais o cobrem: é a tabela
``invertida''.
\begin{align*}
  \min\ & \sum_\ell y_\ell \qquad(\text{ou } \textstyle\sum_\ell f_\ell\,y_\ell \text{ no item b})\\
  \text{s.a.}\ & y_1 + y_6 \ge 1 \ (\text{bairro 1}), \quad y_1 + y_2 \ge 1 \ (2), \quad y_1 + y_3 \ge 1 \ (3), \quad y_2 + y_4 \ge 1 \ (4)\\
  & y_2 + y_3 \ge 1 \ (5), \quad y_3 + y_5 \ge 1 \ (6), \quad y_4 + y_5 \ge 1 \ (7), \quad y_5 + y_6 \ge 1 \ (8)
\end{align*}
\porque ``Pelo menos uma base que cobre o bairro'' é uma soma $\ge1$ sobre os locais que o cobrem. Esse é o modelo
de \emph{cobertura de conjuntos} (set covering).

\passo{Passo 2 --- Solução}
(a) \textbf{3 bases}: L1, L2 e L5. É fácil provar que 2 não bastam: cada local cobre no máximo 3 bairros, e $2\times3=6<8$.
(b) Custo mínimo \textbf{145}, com as mesmas bases ($50+40+55$).

\begin{nota}{Para pensar}
Se a exigência fosse ``cada bairro coberto por \emph{exatamente} uma base'', o modelo viraria \emph{particionamento}
($=1$), visto no exercício A1. Com $\ge1$, coberturas duplas são permitidas: o bairro 2 aqui é coberto por L1 e por L2.
\end{nota}

% ---------------------------------------------------------------------
\exercicio{B4. Escala de enfermeiros}
\enunciado Um hospital precisa de pelo menos o número de enfermeiros da tabela em cada dia da semana. Cada enfermeiro
trabalha 5 dias consecutivos e folga 2 (por exemplo, quem começa na sexta trabalha de sexta a terça). Quantos
enfermeiros contratar, no mínimo?
\begin{center}
\begin{tabular}{lccccccc}
\toprule
Dia & Seg & Ter & Qua & Qui & Sex & Sáb & Dom\\
\midrule
Mínimo & 17 & 13 & 15 & 19 & 14 & 16 & 11\\
\bottomrule
\end{tabular}
\end{center}

\passo{Passo 1 --- A escolha da variável (o ponto principal)}
A tentação é usar ``número de enfermeiros trabalhando no dia $d$''. Isso não funciona, porque não captura a regra dos
5 dias consecutivos. A variável certa é
\[
  x_d \in \mathbb{Z}_{\ge0}: \text{ número de enfermeiros que \emph{começam} a jornada no dia } d .
\]
\porque Sabendo o dia de início, sabemos exatamente em quais dias a pessoa trabalha. A regra da jornada fica
embutida na própria definição da variável.

\passo{Passo 2 --- Modelo}
No dia $d$ trabalha quem começou em $d$ ou nos 4 dias anteriores (contando de forma circular):
\begin{align*}
  \min\ & \sum_d x_d\\
  \text{s.a.}\ & x_{\text{Seg}} + x_{\text{Dom}} + x_{\text{Sáb}} + x_{\text{Sex}} + x_{\text{Qui}} \ge 17 && \text{(segunda)}\\
  & x_{\text{Ter}} + x_{\text{Seg}} + x_{\text{Dom}} + x_{\text{Sáb}} + x_{\text{Sex}} \ge 13 && \text{(terça)}\\
  & \cdots \text{ (uma restrição por dia, sempre o dia e os 4 anteriores)}\\
  & x_d \in \mathbb{Z}_{\ge0}
\end{align*}

\passo{Passo 3 --- Solução e prova de otimalidade}
Ótimo: \textbf{23} enfermeiros. Uma escala: começam 2 na segunda, 6 na terça, 7 na quinta, 3 no sábado e 5 no domingo.

A relaxação linear vale $22{,}33$. Como o número de enfermeiros é inteiro, ele é pelo menos $\lceil22{,}33\rceil=23$, e a
solução acima atinge 23. Ela é ótima sem precisar de mais argumentos.

\begin{nota}{Erro típico}
Arredondar para cima cada valor da relaxação:
\[
  x\approx(1{,}33;\ 5{,}33;\ 0;\ 7{,}33;\ 0;\ 3{,}33;\ 5) \ \longrightarrow\ (2, 6, 0, 8, 0, 4, 5),
\]
o que dá 25 enfermeiros, uma solução viável mas não ótima. Arredondar pode ser muito pior que
resolver o problema inteiro.
\end{nota}

% =====================================================================
\clearpage
\section*{Nível 2 --- Intermediário}
\addcontentsline{toc}{section}{Nível 2 --- Intermediário}

\exercicio{I1. Produção com custo fixo}
\enunciado Uma confecção pode fabricar três produtos (A, B, C), mas cada um exige alugar uma máquina específica, com
custo fixo semanal. Dados por unidade e disponibilidades da semana:
\begin{center}
\begin{tabular}{lcccc}
\toprule
 & A & B & C & Disponível\\
\midrule
Lucro unitário (R\$) & 6 & 4 & 7 & ---\\
Mão de obra (h) & 3 & 2 & 6 & 150\\
Tecido (m$^2$) & 4 & 3 & 4 & 160\\
Aluguel da máquina (R\$) & 200 & 150 & 100 & ---\\
\bottomrule
\end{tabular}
\end{center}
Maximize o lucro semanal (lucro das vendas menos aluguéis).

\passo{Passo 1 --- Variáveis}
$x_i\in\mathbb{Z}_{\ge0}$ (unidades do produto $i$) e $y_i\in\{0,1\}$ (alugar a máquina do produto $i$).

\passo{Passo 2 --- Modelo}
\begin{align*}
  \max\ & 6x_A + 4x_B + 7x_C - 200y_A - 150y_B - 100y_C\\
  \text{s.a.}\ & 3x_A + 2x_B + 6x_C \le 150 && \text{(mão de obra)}\\
  & 4x_A + 3x_B + 4x_C \le 160 && \text{(tecido)}\\
  & x_i \le M_i\,y_i \quad i\in\{A,B,C\} && \text{(só produz com a máquina alugada)}
\end{align*}

\passo{Passo 3 --- Os $M_i$ a partir dos dados}
O maior $x_i$ possível sozinho é limitado pelos dois recursos: $M_i=\min\{150/h_i,\ 160/t_i\}$, arredondado para baixo
porque $x_i$ é inteiro:
\[
  M_A=\min\{50, 40\}=40, \qquad M_B=\min\{75, 53{,}3\}=53, \qquad M_C=\min\{25, 40\}=25 .
\]

\passo{Passo 4 --- Solução e verificação à mão}
Ótimo: \textbf{75}, fabricando só 25 unidades de C. Fabricando um único produto: A dá $6\cdot40-200=40$, B dá
$4\cdot53-150=62$ e C dá $7\cdot25-100=75$. Já qualquer combinação de dois ou mais produtos paga pelo menos
$100+150=250$ de aluguel, e o maior lucro variável que os recursos permitem é 250 (calculado pelo solver), então ela rende no máximo 0.

\begin{nota}{A lição}
Ignorando os aluguéis, o melhor plano é $(30, 0, 10)$, com lucro variável 250. Mas ele usa duas máquinas, e o lucro
real dele é $250-300=-50$. Custo fixo muda completamente a decisão.
\end{nota}

% ---------------------------------------------------------------------
\exercicio{I2. Restrições ``ou-ou'' e ``se-então''}
\enunciado Uma empresa fabrica dois produtos ($x_1$ e $x_2$ unidades inteiras, no máximo 8 de cada) e quer maximizar
$3x_1+4x_2$. A produção será feita em \emph{uma} de duas plantas: na planta 1 vale $2x_1+3x_2\le18$; na planta 2 vale
$4x_1+2x_2\le20$. Além disso, se o produto 1 for fabricado, a embalagem só comporta até 3 unidades do produto 2.

\passo{Passo 1 --- Ou-ou: escolher a planta}
$w\in\{0,1\}$, com $w=0$ para a planta 1 e $w=1$ para a planta 2:
\begin{align*}
  2x_1 + 3x_2 &\le 18 + M_1\,w\\
  4x_1 + 2x_2 &\le 20 + M_2\,(1-w)
\end{align*}
Com $w=0$, só a primeira fica ativa; com $w=1$, só a segunda. Os $M$ vêm dos limites $x_1,x_2\le8$: o maior valor
possível do lado esquerdo menos o lado direito. $M_1=2\cdot8+3\cdot8-18=22$ e $M_2=4\cdot8+2\cdot8-20=28$.

\passo{Passo 2 --- Se-então: a regra da embalagem}
``Se $x_1>0$, então $x_2\le3$'' liga uma condição sobre uma variável inteira a uma restrição. Use uma binária $z$ como
intermediária:
\begin{align*}
  x_1 &\le 8\,z && \text{($x_1>0$ obriga $z=1$)}\\
  x_2 &\le 3 + 5\,(1-z) && \text{($z=1$ ativa $x_2\le3$; com $z=0$ fica $x_2\le8$)}
\end{align*}
\porque Não dá para escrever ``$x_1>0$'' diretamente. Quebramos a implicação em duas partes:
$x_1>0\Rightarrow z=1$ e $z=1\Rightarrow x_2\le3$. O 5 é o big-M desta restrição: $8-3$.

\passo{Passo 3 --- Solução}
Ótimo: \textbf{32}, com $x=(0, 8)$ na planta 2. Sem a regra da embalagem o ótimo seria 35, com $x=(1, 8)$ na planta 2, mas
aí o produto 1 é fabricado junto com 8 unidades do produto 2, o que a regra proíbe. Na planta 1 o melhor seria 26, com
$x=(6,2)$. Os valores foram conferidos por enumeração de todos os 81 pares $(x_1,x_2)$.

% ---------------------------------------------------------------------
\exercicio{I3. Corte unidimensional (padrões de corte)}
\enunciado Uma serralheria corta barras de 10\,m para atender pedidos de 14 peças de 3\,m, 10 peças de 4\,m e
7 peças de 5\,m. Quantas barras usar, no mínimo?

\passo{Passo 1 --- A ideia: decidir \emph{padrões}, não peças}
Um \emph{padrão de corte} é uma forma de cortar uma barra. Listamos os padrões \emph{maximais}, aqueles em que não cabe
mais nenhuma peça na sobra:
\begin{center}
\begin{tabular}{ccccc}
\toprule
Padrão & Peças de 3\,m & Peças de 4\,m & Peças de 5\,m & Sobra (m)\\
\midrule
$p_1$ & 3 & 0 & 0 & 1\\
$p_2$ & 2 & 1 & 0 & 0\\
$p_3$ & 1 & 0 & 1 & 2\\
$p_4$ & 0 & 2 & 0 & 2\\
$p_5$ & 0 & 1 & 1 & 1\\
$p_6$ & 0 & 0 & 2 & 0\\
\bottomrule
\end{tabular}
\end{center}
\porque A alternativa (``a peça $k$ vai para a barra $b$'') cria muitas variáveis e muitas soluções equivalentes, porque
as barras são idênticas. Com padrões, a pergunta vira apenas ``quantas barras corto de cada jeito?''.

\passo{Passo 2 --- Modelo}
$x_p\in\mathbb{Z}_{\ge0}$: número de barras cortadas no padrão $p$; $a_{kp}$ é o número de peças do tamanho $k$ no padrão $p$.
\begin{align*}
  \min\ & \sum_p x_p\\
  \text{s.a.}\ & 3x_1 + 2x_2 + x_3 \ge 14 && \text{(peças de 3\,m)}\\
  & x_2 + 2x_4 + x_5 \ge 10 && \text{(peças de 4\,m)}\\
  & x_3 + x_5 + 2x_6 \ge 7 && \text{(peças de 5\,m)}
\end{align*}
O $\ge$ permite peças a mais, que viram estoque. Essa é a formulação de Gilmore e Gomory. Em problemas grandes, os padrões
não são listados todos de antemão: eles são gerados conforme a necessidade (geração de colunas).

\passo{Passo 3 --- Solução e prova}
Ótimo: \textbf{12 barras}: 7 no padrão $p_2$, 3 no $p_6$, 1 no $p_4$ e 1 no $p_5$. O comprimento total pedido é
$42+40+35=117$\,m, então são necessárias pelo menos $\lceil117/10\rceil=12$ barras, e a solução atinge esse limite.

% ---------------------------------------------------------------------
\exercicio{I4. Desconto por quantidade (custo linear por partes)}
\enunciado Uma fábrica precisa comprar 420\,kg de uma resina. O fornecedor 1 dá desconto incremental: os primeiros 100\,kg
custam R\$\,10/kg, do 101.º ao 300.º kg custam R\$\,8/kg e do 301.º ao 500.º kg custam R\$\,6/kg. O fornecedor 2 cobra
R\$\,7,50/kg, vende no máximo 250\,kg e cobra frete fixo de R\$\,300 se for usado. Minimize o custo da compra.

\passo{Passo 1 --- Dividir a compra do fornecedor 1 em faixas}
$s_1\in[0,100]$, $s_2\in[0,200]$, $s_3\in[0,200]$: quantidade comprada em cada faixa de preço. A compra do fornecedor 1 é
$s_1+s_2+s_3$ e o custo é $10s_1+8s_2+6s_3$.

\passo{Passo 2 --- Por que é preciso binárias aqui}
O preço \emph{cai} com a quantidade (custo côncavo). Sem nenhuma restrição extra, o solver, que minimiza, compraria direto na
faixa mais barata ($s_3$) sem ``passar'' pelas faixas caras, o que é impossível na vida real. Precisamos forçar a ordem:
só se compra na faixa $k+1$ se a faixa $k$ estiver cheia. Com $z_k\in\{0,1\}$ indicando que a faixa $k$ está em uso:
\begin{align*}
  s_k &\le \ell_k\,z_k && k=1,2,3 \quad(\ell = 100, 200, 200)\\
  s_k &\ge \ell_k\,z_{k+1} && k=1,2 \quad\text{(usar a faixa seguinte exige esta cheia)}
\end{align*}
\porque Se o preço \emph{subisse} com a quantidade (custo convexo), as binárias seriam desnecessárias: o solver encheria as
faixas baratas primeiro por conta própria. Custo côncavo num problema de mínimo é o caso que exige binárias.

\passo{Passo 3 --- Modelo completo}
\begin{align*}
  \min\ & 10s_1 + 8s_2 + 6s_3 + 7{,}5\,q + 300\,u\\
  \text{s.a.}\ & s_1 + s_2 + s_3 + q = 420\\
  & q \le 250\,u && \text{(fornecedor 2 só com frete pago)}\\
  & \text{restrições de faixa do passo 2}, \quad z_k, u\in\{0,1\}
\end{align*}

\passo{Passo 4 --- Solução}
Ótimo: \textbf{R\$\,3320}, comprando tudo do fornecedor 1: $100\cdot10+200\cdot8+120\cdot6=1000+1600+720$. Usar o fornecedor 2 ao
máximo custaria $250\cdot7{,}5+300=2175$, mais 170\,kg do fornecedor 1 por $1000+560=1560$, total 3735.

\begin{nota}{Erro típico}
Esquecer as restrições de faixa. O modelo devolve R\$\,3000, com $s=(20, 200, 200)$: compra 200\,kg a R\$\,6 tendo comprado
só 20\,kg da primeira faixa, um preço que o fornecedor nunca daria.
\end{nota}

% =====================================================================
\clearpage
\section*{Nível 3 --- Avançado}
\addcontentsline{toc}{section}{Nível 3 --- Avançado}

\exercicio{A1. Escala de tripulações: particionamento \texorpdfstring{$\times$}{x} cobertura}
\enunciado Uma companhia aérea precisa atribuir tripulações a 6 voos. O setor de planejamento gerou 10 \emph{jornadas}
viáveis (sequências de voos que uma tripulação consegue fazer em um dia, respeitando as regras trabalhistas), cada uma
com seu custo. Cada voo precisa de exatamente uma tripulação trabalhando.
\begin{center}
\begin{tabular}{lcccccccccc}
\toprule
Jornada & J1 & J2 & J3 & J4 & J5 & J6 & J7 & J8 & J9 & J10\\
\midrule
Voos & 1,2 & 3,4 & 5,6 & 1,3,5 & 2,4,6 & 1,4 & 2,3,6 & 5 & 1,2,3,4 & 4,5,6\\
Custo & 5 & 6 & 6 & 9 & 8 & 4 & 7 & 3 & 7 & 5\\
\bottomrule
\end{tabular}
\end{center}
(a) Modele o problema. (b) O que muda se for permitido que uma tripulação viaje como passageira em um voo já coberto
(\emph{deadhead})?

\passo{Passo 1 --- Variáveis e modelo de particionamento}
$x_j\in\{0,1\}$: a jornada $j$ é escolhida. Para cada voo, a soma das jornadas que o contêm deve ser \emph{exatamente} 1:
\begin{align*}
  \min\ & \sum_j c_j\,x_j\\
  \text{s.a.}\ & x_1 + x_4 + x_6 + x_9 = 1 && \text{(voo 1)}\\
  & x_1 + x_5 + x_7 + x_9 = 1 && \text{(voo 2)}\\
  & x_2 + x_4 + x_7 + x_9 = 1 && \text{(voo 3)}\\
  & x_2 + x_5 + x_6 + x_9 + x_{10} = 1 && \text{(voo 4)}\\
  & x_3 + x_4 + x_8 + x_{10} = 1 && \text{(voo 5)}\\
  & x_3 + x_5 + x_7 + x_{10} = 1 && \text{(voo 6)}
\end{align*}
\porque Toda a complexidade das regras trabalhistas ficou na \emph{geração} das jornadas. O modelo de otimização só
escolhe uma combinação delas. É a mesma ideia dos padrões de corte (I3): enumerar as opções viáveis e escolher entre elas.

\passo{Passo 2 --- Solução (a)}
Ótimo do particionamento: \textbf{13}, com J9 (voos 1 a 4) e J3 (voos 5 e 6).

\passo{Passo 3 --- Com deadhead (b): cobertura}
Se uma tripulação pode viajar como passageira, um voo pode aparecer em mais de uma jornada escolhida. As restrições
passam a ser $\ge1$ (cobertura). O ótimo cai para \textbf{12}, com J9 e J10: o voo 4 é coberto pelas duas, e uma das
tripulações viaja nele como passageira.

\begin{nota}{A lição}
A diferença entre $=1$ e $\ge1$ não é detalhe de notação: ela representa uma regra operacional. A cobertura é mais
fácil de resolver (tem mais soluções viáveis) e dá um limitante inferior para o particionamento.
\end{nota}

% ---------------------------------------------------------------------
\exercicio{A2. Programação de tarefas em job shop}
\enunciado Três ordens de produção passam por três máquinas, cada uma com seu próprio roteiro (a ordem das máquinas é
fixa para cada ordem de produção). Cada máquina faz uma operação por vez, sem interrupção. Minimize o instante em que tudo
termina (makespan).
\begin{center}
\begin{tabular}{ll}
\toprule
Ordem & Roteiro (máquina, tempo)\\
\midrule
1 & M1 (3) $\to$ M2 (2) $\to$ M3 (2)\\
2 & M1 (2) $\to$ M3 (1) $\to$ M2 (4)\\
3 & M2 (4) $\to$ M3 (3)\\
\bottomrule
\end{tabular}
\end{center}

\passo{Passo 1 --- Variáveis}
$S_{jo}\ge0$: início da operação $o$ da ordem $j$; $C_{\max}$: makespan; e, para cada par de operações que usam a mesma
máquina, uma binária $y$ que escolhe qual vem primeiro.

\passo{Passo 2 --- Dois tipos de restrição}
\begin{enumerate}[label=(\roman*)]
  \item \textbf{Roteiro} (dentro de uma ordem): a operação seguinte só começa quando a anterior termina, por exemplo
        $S_{12}\ge S_{11}+3$ e $S_{13}\ge S_{12}+2$. A ordem é fixa, então não há escolha nem binária.
  \item \textbf{Máquina} (entre ordens): duas operações na mesma máquina não se sobrepõem. É uma disjunção, como no
        exercício 3.1 dos Desafios. Para as operações $(j,o)$ e $(k,q)$ na mesma máquina:
        \[
          S_{jo} + p_{jo} \le S_{kq} + H\,(1-y), \qquad S_{kq} + p_{kq} \le S_{jo} + H\,y .
        \]
\end{enumerate}
Mais $C_{\max}\ge S_{j,\text{última}}+p_{j,\text{última}}$ para cada ordem, e $\min C_{\max}$.

\passo{Passo 3 --- Quantas binárias e qual $M$?}
Pares na mesma máquina: M1 tem as ordens 1 e 2 (1 par); M2 tem as ordens 1, 2 e 3 (3 pares); M3 tem as ordens 1, 2 e 3 (3 pares).
São 7 binárias. Para o $M$: executar todas as operações em sequência, sem paralelismo, termina em
$H=\sum p=7+7+7=21$, então nenhuma solução precisa passar de 21 e $M=H=21$ basta.

\passo{Passo 4 --- Solução}
Ótimo: \textbf{11}.
\begin{center}
\begin{tabular}{lll}
\toprule
Máquina & Sequência (intervalos) & \\
\midrule
M1 & ordem 2 [0, 2], ordem 1 [2, 5] & \\
M2 & ordem 3 [0, 4], ordem 1 [5, 7], ordem 2 [7, 11] & \\
M3 & ordem 2 [2, 3], ordem 3 [4, 7], ordem 1 [7, 9] & \\
\bottomrule
\end{tabular}
\end{center}
Limitante simples: a máquina M2 tem $2+4+4=10$ horas de trabalho, então o makespan é pelo menos 10. O solver prova que
10 não é atingível por causa dos roteiros.

% ---------------------------------------------------------------------
\exercicio{A3. Roteamento de veículos com capacidade}
\enunciado Um depósito em $(0,0)$ atende 6 clientes com veículos de capacidade 15 (há 3 veículos disponíveis). A distância
é a de Manhattan. Cada cliente é visitado por exatamente um veículo; cada veículo sai do depósito e volta a ele, e a carga
de uma rota não pode passar de 15. Minimize a distância total.
\begin{center}
\begin{tabular}{lcccccc}
\toprule
Cliente & 1 & 2 & 3 & 4 & 5 & 6\\
\midrule
Posição & (2, 4) & (5, 3) & (6, $-1$) & ($-3$, 3) & ($-4$, $-2$) & (1, $-5$)\\
Demanda $q_i$ & 4 & 6 & 5 & 7 & 4 & 6\\
\bottomrule
\end{tabular}
\end{center}

\passo{Passo 1 --- Variáveis}
$x_{ij}\in\{0,1\}$ (um veículo vai de $i$ a $j$); $u_i$: carga acumulada na rota até o cliente $i$, inclusive, com
$q_i\le u_i\le15$; e $K\in\{0,\dots,3\}$, o número de veículos usados.

\passo{Passo 2 --- Modelo}
\begin{align*}
  \min\ & \sum_{i\ne j} d_{ij}\,x_{ij}\\
  \text{s.a.}\ & \sum_{j} x_{ij} = 1, \quad \sum_{j} x_{ji} = 1 && \forall\,\text{cliente } i\\
  & \sum_{j} x_{0j} = K, \quad \sum_{j} x_{j0} = K && \text{(saem e voltam $K$ veículos)}\\
  & u_j \ge u_i + q_j - 15\,(1 - x_{ij}) && \forall\,\text{clientes } i\ne j\\
  & q_i \le u_i \le 15
\end{align*}
\porque A restrição em $u$ é o MTZ do caixeiro (exercício 1.2 dos Desafios), com a carga no lugar da posição. Se $x_{ij}=1$, a carga
em $j$ é a carga em $i$ mais a demanda de $j$. Ela faz dois trabalhos de uma vez: impede subciclos sem o depósito (a carga
teria de crescer para sempre) e, com $u_i\le15$, garante a capacidade do veículo. O big-M é a própria capacidade 15: com
$x_{ij}=0$ fica $u_j\ge u_i+q_j-15$, sempre verdade porque $u_i\le15$ e $u_j\ge q_j$.

\passo{Passo 3 --- Solução}
A demanda total é 32 e cada veículo leva 15, então são necessários pelo menos $\lceil32/15\rceil=3$ veículos. Ótimo:
\textbf{52}, com as rotas
\begin{center}
\begin{tabular}{lcc}
\toprule
Rota & Carga & Distância\\
\midrule
$0\to3\to2\to1\to0$ & $5+6+4=15$ & $7+5+4+6=22$\\
$0\to5\to4\to0$ & $4+7=11$ & $6+6+6=18$\\
$0\to6\to0$ & 6 & $6+6=12$\\
\bottomrule
\end{tabular}
\end{center}

\begin{nota}{Erro típico}
Usar só a capacidade ``por rota'' sem uma variável de carga. Sem $u_i$ não há como saber quais clientes estão na mesma rota,
e o modelo aceitaria subciclos que nunca passam pelo depósito.
\end{nota}

% =====================================================================
\clearpage
\section*{Resumo: qual técnica aparece em cada exercício}
\addcontentsline{toc}{section}{Resumo: qual técnica aparece em cada exercício}
\begin{center}
\begin{tabular}{>{\raggedright}p{3.2cm}>{\raggedright}p{6.2cm}>{\raggedright\arraybackslash}p{4.8cm}}
\toprule
Exercício & Técnica principal & Ideia-chave\\
\midrule
B1 Projetos & mochila + regras lógicas & $\le$ para ``se'', $=$ para ``os dois ou nenhum''\\
B2 Designação & designação & relaxação já é inteira\\
B3 Ambulâncias & cobertura de conjuntos & soma $\ge1$ por bairro\\
B4 Enfermeiros & variável ``quem começa no dia $d$'' & $\lceil$relaxação$\rceil$ prova o ótimo\\
I1 Custo fixo & $x\le M_i\,y$ com $M_i$ dos recursos & custo fixo muda a decisão\\
I2 Plantas & ou-ou e se-então & binária intermediária $z$\\
I3 Corte & padrões de corte (Gilmore--Gomory) & decidir padrões, não peças\\
I4 Desconto & linear por partes côncavo & faixas em ordem exigem binárias\\
A1 Tripulações & particionamento $\times$ cobertura & $=1$ ou $\ge1$ é regra operacional\\
A2 Job shop & roteiro + disjunções por máquina & $M=\sum p$\\
A3 Roteamento & MTZ com carga & $u$ elimina subciclos e limita carga\\
\bottomrule
\end{tabular}
\end{center}

\section*{Referência}
Arenales, M.; Armentano, V.; Morabito, R.; Yanasse, H. \emph{Pesquisa Operacional}. Rio de Janeiro: Elsevier/Campus, 2007.

\end{document}
